\begin{partdivider}{III}{LE TRAVAIL ET L'INTELLIGENCE}\end{partdivider}

\bookchap{7}{Le travail ne disparaît pas : il se transforme}

Dès que l'on évoque l'intelligence artificielle et le travail, un débat sempiternel ressurgit, souvent avec une simplicité déconcertante : l'IA va-t-elle supprimer les emplois ? Ce débat oppose généralement deux camps irréconciliables. D'un côté, les apocalyptiques annoncent le chômage de masse, la fin du travail et l'obsolescence de l'humain. De l'autre, les optimistes rassurent en affirmant que, comme chaque révolution industrielle, celle-ci créera plus d'emplois qu'elle n'en détruira.

Ce débat, tel qu'il est formulé, est impuissant à éclairer notre avenir. Il échoue parce qu'il raisonne à l'échelle trop macroscopique de «l'emploi» en général, et parce qu'il ignore la nature exacte de ce que l'IA modifie.

Pour comprendre ce qui se joue, il faut introduire une distinction fondamentale, déjà effleurée dans la première partie de ce livre, mais qui prend ici tout son sens :

\begin{aiquote}[s]
\qbold L'IA n'automatise pas des métiers. Elle automatise des tâches.
\end{aiquote}

Un métier n'est jamais un bloc monolithique. Il est un faisceau de tâches diverses, certaines répétitives, d'autres créatives, certaines d'exécution, d'autres de jugement. Prenons l'exemple d'un radiologue. Son métier inclut l'analyse d'images médicales (une tâche de reconnaissance visuelle de plus en plus automatisable), mais aussi l'histoire du patient, la recommandation thérapeutique, l'annonce du diagnostic et le réconfort \textbackslash\{\}1

Les travaux récents le confirment : une étude menée par les chercheurs d'OpenAI en 2023 a ainsi montré que la majorité des professions humaines ont au moins une fraction de leurs tâches susceptibles d'être accélérées ou automatisées par les grands modèles de langage. L'impact ne se limite plus aux travaux manuels ou administratifs de routine ; il touche le cœur du travail intellectuel.

C'est un fait établi : lorsqu'une tâche est automatisée, le métier ne disparaît pas nécessairement ; il se recompose. Les tâches restantes changent de proportion, d'importance et de nature. Le temps libéré par l'automatisation d'une tâche peut être réalloué à d'autres tâches, souvent plus qualitatives, ou bien il peut conduire à une réduction de l'effectif nécessaire pour accomplir le même volume de travail. C'est cette recomposition qui est le véritable enjeu, bien plus que la destruction pure et simple. Et elle est déjà à l'œuvre.

Considérons l'histoire récente des distributeurs automatiques de billets (DAB). Lors de leur apparition, beaucoup craignirent la disparition des guichetiers de banque. Or, le nombre d'employés de banque n'a pas chuté de manière spectaculaire dans de nombreux pays. Ce qui a changé, c'est la nature de leur travail : les guichetiers sont devenus des conseillers. La tâche de distribution de billets, purement d'exécution, a été automatisée, libérant du temps pour des tâches de conseil, de relation et de vente. Le métier s'est rehaussé en complexité.

L'actualité récente offre des illustrations frappantes de cette recomposition. La grève historique des scénaristes et des acteurs à Hollywood en 2023 a offert un aperçu saisissant de ce débat : la défense ne portait pas seulement sur la crainte d'une automatisation de l'écriture de scripts par l'IA, mais sur l'utilisation non consentie des images et des voix pour créer des répliques numériques éternelles, modifiant à jamais le contrat de travail et la propriété de l'empreinte humaine.

L'hypothèse qui se dessine est la suivante : avec l'IA, ce mouvement de recomposition va s'accélérer et s'étendre à des sphères qui, jusqu'ici, étaient épargnées parce qu'elles relevaient du travail intellectuel. Les tâches de recherche, de synthèse, de rédaction standard, de codage répétitif ou d'analyse de données vont être progressivement absorbées par la machine. Le travail humain se déplacera vers ce que la machine fait moins bien : l'interaction humaine, la supervision, la gestion des cas limites, la décision dans l'incertitude et l'imagination de nouveaux usages.

Ce déplacement crée également de nouvelles tâches, qui n'existaient pas ou qui étaient marginales. Qui évalue la qualité d'une réponse générée par une IA ? Qui conçoit les prompts pertinents pour orienter la machine ? Qui audite les biais du système ? Qui assure la transition et la formation des équipes ? Ces nouvelles tâches font émerger de nouvelles compétences et, parfois, de nouvelles professions.

Prenons un exemple concret dans le développement logiciel. L'arrivée des copilotes de code n'a pas éliminé le besoin de développeurs. En revanche, elle modifie leur quotidien. Le développeur passe moins de temps à écrire des lignes de code syntaxiquement correctes, et davantage de temps à spécifier l'architecture, à vérifier la logique du code proposé, à optimiser les interactions et à penser aux tests de sécurité. Son métier glisse d'exécutant-codeur vers architecte-superviseur.

Il faut cependant regarder cette transformation en face, sans édulcorer ses conséquences sociales. Si le métier se recompose, c'est aussi parce que la productivité augmente. Si une tâche qui nécessitait cinq personnes n'en demande plus que deux assistées par l'IA, alors, à volume de production constant, trois postes sont potentiellement excédentaires. La recomposition des tâches ne garantit pas la conservation des effectifs.

De plus, la recomposition exige une adaptabilité considérable de la part des travailleurs. Passer de l'exécution d'une tâche technique à la supervision d'un système intelligent demande une reconceptualisation de son identité professionnelle, et souvent une formation. Or, les systèmes sociaux et les individus ont une inertie bien plus grande que la technologie. C'est dans cet écart — entre la vitesse de la recomposition technique et la lenteur de l'adaptation humaine — que réside le risque principal de fracture sociale et professionnelle.

La question n'est donc pas de savoir si le travail va disparaître. Il ne disparaîtra pas. La question est de savoir \emph{comment} il se transforme, \emph{à quel rythme}, et si nos institutions, nos entreprises et nos systèmes de formation permettent d'accompagner cette transformation plutôt que de la subir.

Pour mieux en saisir la dynamique, il nous faut maintenant zoomer au niveau de l'individu. Que se passe-t-il lorsqu'un professionnel s'adjoint les services d'une IA ?

\bookchap{8}{Le professionnel augmenté}

Imaginons un instant un monde du travail qui aurait franchi le pas de l'intégration massive de l'IA. Dans ce monde, chaque professionnel disposerait, non pas d'un simple logiciel, mais d'un ensemble d'agents spécialisés qui l'accompagneraient dans son activité quotidienne.

Un entrepreneur aurait un agent stratégique pour analyser les marchés, un agent juridique pour vérifier les contrats, et un agent de communication pour gérer sa correspondance. Un médecin disposerait d'un copilote de diagnostic croisant les données du patient avec la littérature médicale mondiale, et d'un agent administratif pour préremplir ses comptes rendus. Un artisan pourrait s'appuyer sur un agent de design 3D pour visualiser ses créations, et sur un agent logistique pour optimiser ses approvisionnements.

On observe que l'une des conséquences les plus immédiates de cette configuration est l'explosion potentielle de la productivité individuelle. Ce qui nécessitait hier une équipe de cinq personnes — un chercheur, un analyste, un rédacteur, un graphiste et un correcteur — pourrait être accompli par un seul individu, augmenté de cinq agents IA.

Ce scénario n'est pas une vue de l'esprit. Des expériences récentes le valorent déjà. Des développeurs utilisant des copilotes de code rapportent des gains de productivité de 30 à 50~\% sur certaines tâches de programmation. Des consultants produisent en quelques heures des rapports qui leur prenaient des jours. Des créateurs de contenu génèrent des dizaines de variantes d'un même texte pour tester différents angles.

La figure du professionnel augmenté est séduisante. Elle dessine l'image d'un travailleur libéré des tâches les plus fastidieuses, recentré sur ce qui lui est propre : la décision, la stratégie, la créativité et la relation. Elle porte en elle la promesse d'une démocratisation de l'entrepreneuriat : si l'on peut à soi seul accomplir le travail d'une petite équipe, la création d'entreprise en est grandement facilitée.

Mais cette même figure soulève une question difficile, qui touche au cœur même de notre organisation économique :

\begin{aiquote}
Que devient la valeur de l'expertise humaine lorsque tout le monde est augmenté ?
\end{aiquote}

Il faut ici distinguer deux situations.

Dans la première, l'augmentation est asymétrique. Seuls certains professionnels ont accès à l'IA. Dans ce cas, l'augmentation est un avantage compétitif formidable. Le professionnel augmenté surpasse massivement ses concurrents non augmentés. Sa valeur individuelle sur le marché augmente.

Dans la seconde situation, l'augmentation se démocratise. Tout le monde, ou presque, dispose de ces agents. L'augmentation devient la norme, le ticket d'entrée minimum pour participer au marché. C'est le scénario le plus probable à terme si la technologie continue de se répandre et de se banaliser.

L'hypothèse est alors la suivante : lorsque l'augmentation, elle cesse d'être un avantage différentiel pour redevenir un standard. La barre de ce qui est attendu se rehausse. Ce qui était remarquable devient normal. Et la valeur économique se déplace à nouveau.

C'est un phénomène bien connu en économie de l'innovation, parfois appelé l'«effet Reine rouge» (en référence à \emph{Alice au pays des merveilles} : il faut courir très vite juste pour rester sur place). Si tous les avocats augmentés peuvent synthétiser la jurisprudence en une heure, le faire ne constitue plus un avantage concurrentiel ; c'est le minimum requis. La valeur se déplace vers d'autres dimensions du métier : la capacité à trouver les bons litiges, à négocier, à convaincre un juge ou à accompagner un client angoissé. De plus, l'augmentation généralisée pose le problème de l'évaluation. Si un junior augmenté peut produire un rapport d'une qualité apparente équivalente à celle d'un senior, comment l'employeur évalue-t-il la véritable compétence ? Le junior a produit un résultat correct, mais il n'a pas forcément la maturité pour en saisir les limites ou pour l'adapter à un contexte non prévu par la machine. L'illusion de compétence, rendue possible par l'IA, risque de brouiller les signaux de qualité sur le marché du travail.

Enfin, il faut considérer un risque plus subtil, que l'on pourrait appeler la

\emph{dégradation de l'expertise par automatisation de l'effort}. L'expertise humaine se construit souvent par la confrontation répétée à la difficulté. C'est en cherchant la solution à un problème de code récalcitrant que le développeur acquiert une intuition architecturale. C'est en rédigeant péniblement un argumentaire qu'un avocat en intériorise les détours. Si l'IA supprime l'effort, elle supprime peut-être avec lui le terreau de l'expertise profonde.

\begin{aiquote}
Si nous n'avons plus à lutter pour produire, que savons-nous encore de ce que nous produisons ?
\end{aiquote}

Le professionnel augmenté est donc une réalité puissante et ambiguë. Il gagne en puissance d'exécution, mais risque de perdre en compréhension intime de son propre travail. Il voit sa productivité exploser, mais se trouve confronté à une nouvelle exigence : il doit prouver sa valeur sur des dimensions (jugement, créativité, éthique) que la machine ne maîtrise pas.

Cette tension entre augmentation et érosion de l'expertise nous conduit à examiner un phénomène qui pourrait modifier profondément notre rapport au savoir : la possibilité, pour chacun, de devenir instantanément un «expert» par procuration.

\bookchap{9}{Quand tout le monde peut devenir expert}

Pendant des siècles, l'expertise a été un rempart. L'expert — le médecin, l'avocat, l'ingénieur, le professeur — occupait une position privilégiée parce qu'il détenait un savoir inaccessible au commun. Son autorité reposait sur la rareté de ses connaissances et sur le temps long de son acquisition.

L'IA vient frapper ce rempart d'un coup inattendu. Non pas en rendant l'expertise inutile, mais en la rendant \emph{consultable} par n'importe qui, en temps réel. Imaginons une situation banale. Une personne sans aucune formation juridique doit rédiger un contrat de bail pour louer son appartement. Autrefois, elle aurait dû consulter un notaire ou un avocat, ou se débrouiller avec un modèle standard qu'elle ne comprenait pas vraiment. Aujourd'hui, elle peut demander à une IA de générer un contrat de bail adapté à la législation de son pays, en y insérant des clauses spécifiques de son choix. En quelques minutes, elle obtient un document d'une qualité formelle souvent supérieure à ce qu'elle aurait pu rédiger seule.

De même, un étudiant qui bute sur un théorème de physique quantique peut obtenir en quelques secondes une explication pédagogique, adaptée à son niveau, avec des analogies et des exemples. Un non-développeur peut créer une petite application fonctionnelle pour gérer ses stocks.

On observe que l'IA permet à un non-expert de produire un résultat qui a l'apparence de l'expertise. Elle fait baisser drastiquement le coût d'accès au savoir spécialisé. La frontière entre celui qui sait et celui qui ne sait pas, si elle ne s'efface pas, devient en tout cas beaucoup plus poreuse.

Cette démocratisation de l'expertise est, sur bien des plans, une bonne nouvelle. Elle peut réduire les asymétries d'information qui, dans de nombreux domaines (santé, droit, finance), plaçaient le citoyen ordinaire en position de faiblesse face aux professionnels. Elle peut permettre à des populations dépourvues d'infrastructures éducatives d'accéder à un niveau de conseil et d'explication inespéré. En Afrique ou dans certaines zones rurales du monde, le fait de pouvoir interroger un système qui dispense des conseils médicaux de premier recours ou des pratiques agricoles adaptées n'est pas un gadget : c'est potentiellement une avancée vitale.

Toutefois, cette démocratisation comporte des risques majeurs, que l'on commencerait à peine à mesurer.

Le premier risque est \textbf{l'illusion de compétence}. Disposer d'une réponse correcte ne signifie pas qu'on la comprend. Pouvoir générer un contrat ne confère pas la capacité de juger si une clause est abusive ou contestable dans un contexte spécifique. L'illusion de compétence est ce biais cognitif qui nous fait croire que nous maîtrisons un sujet parce que nous avons obtenu un bon résultat avec l'aide d'un outil. C'est l'effet Dunning-Kruger amplifié par la machine : moins on sait, moins on est capable d'évaluer les limites de ce que la machine nous donne, et plus on a l'illusion de savoir.

Le deuxième risque est la \textbf{dépendance sans filet de sécurité}. Si l'IA se trompe — et elle se trompe, car elle «hallucine» des informations plausibles mais fausses —, l'utilisateur non expert n'a souvent aucun moyen de s'en apercevoir. L'expert humain, confronté à une erreur, dispose d'un réseau de intuitions, de connaissances croisées et de savoir-faire qui lui fait souvent sentir l'incohérence avant même de la démontrer. Le non-expert augmenté, lui, est démuni. Il valide la réponse par défaut, parce qu'elle a l'air cohérente et qu'il n'a pas de critère de rejet. C'est le paradoxe du pilote automatique : tant que tout va bien, on va plus vite ; dès que le système échoue, on a perdu les réflexes pour reprendre la main.

Le troisième risque est \textbf{l'érosion du savoir-faire intermédiaire}. Dans toute chaîne de valeur, il existe des métiers de demi-expertise : le technicien, le parajuriste, l'infirmière, le développeur junior. Ce sont souvent eux qui font le lien entre l'expert et le terrain, qui assurent l'exécution, qui détectent les anomalies. Si l'IA permet à l'expert de court-circuiter ces intermédiaires en produisant directement l'exécution, à qui l'expert délègue-t-il la vérification de terrain ? Et comment les jeunes professionnels acquièrent-ils l'expérience si les tâches de base qui formaient leur apprentissage sont automatisées ?

L'interprétation est alors la suivante : l'IA ne supprime pas le besoin d'experts. Elle modifie la nature de leur fonction. L'expertise se déplace de la capacité à \emph{produire} une réponse vers la capacité à \emph{auditer} une réponse, à en évaluer la pertinence dans un contexte réel, et à en assumer la responsabilité.

L'expert de demain ne sera peut-être plus celui qui sait tout, mais celui qui sait détecter les limites de ce que la machine propose, qui connaît les angles morts du système et qui peut intervenir quand le modèle, par construction, échoue. C'est une forme d'expertise plus méta, plus critique, plus situationnelle.

Mais comment former de tels experts ? Comment transmettre ce jugement critique si l'on n'est plus passé par le lent apprentissage de l'exécution ? C'est tout le défi qui se pose à notre système éducatif.

Mais cette illusion de compétence à l'échelle individuelle possède un pendant systémique beaucoup plus redoutable. Si l'utilisateur croit naïvement en l'infaillibilité ou en la neutralité de la machine, il risque de ne pas voir qui, de l'autre côté de l'écran, oriente réellement cette machine, sélectionne ses données d'apprentissage et en tire les bénéfices. L'expertise qui s'érode laisse la place à un nouveau pouvoir, invisible et concentré. C'est ce pouvoir qu'il nous faut maintenant interroger.

\bookchap{10}{L'école après Google}

De toutes les institutions que l'IA pourrait bouleverser, l'école est peut-être celle où le choc est le plus existentiel. Non pas parce que les enseignants vont disparaître, mais parce que la raison d'être même de l'école est interrogée.

Pendant des siècles, l'école a eu pour fonction principale la transmission d'un savoir accumulé. Mémoriser, restituer, appliquer des règles : sur ces piliers, l'institution scolaire a bâti ses rituels, ses évaluations et ses hiérarchies. Le contrôle des connaissances supposait que l'accès à ces connaissances était limité. Le professeur était le dépositaire, le manuel le médiateur, l'examen le vérificateur.

Puis Internet est arrivé, et avec lui, l'accès quasi illimité à l'information. L'école a commencé à vaciller. À quoi bon mémoriser des dates ou des formules si elles sont à portée de clic ? Le système a résisté, arguant que mémoriser restait indispensable pour penser, et que l'information n'était pas le savoir.

Avec l'IA générative, le séisme est d'une autre ampleur. Ce n'est plus seulement l'accès à l'information qui est facile, c'est la production de raisonnement à partir de cette information. Si une IA peut, en quelques secondes, répondre à une question de philosophie, résoudre une équation, traduire un texte, résumer un chapitre ou programmer un algorithme, alors l'évaluation classique des compétences cognitives de base perd une grande partie de son sens.

\begin{aiquote}
Que devons-nous encore apprendre, et comment, si la machine peut faire nos devoirs ?
\end{aiquote}

La tentation est grande d'interdire. Interdire l'IA à l'école, comme on a un temps tenté d'interdire la calculatrice ou le smartphone. Cette tentation est compréhensible, mais elle est probablement vaine. On n'interdit pas une technologie qui s'infuse dans tous les pans de la société. De plus, elle est contre-productive : elle prive les élèves de l'outil avec lequel ils devront travailler, et elle laisse l'école s'enfermer dans un anachronisme de plus en plus déconnecté du réel.

L'hypothèse qui se dessine est la suivante : le défi de l'école n'est pas de protéger le savoir de l'IA, mais de former des esprits capables de penser \emph{avec} l'IA, \emph{contre} l'IA et

\emph{au-delà} de l'IA.

Cela suppose un profond déplacement du centre de gravité de l'enseignement. Un déplacement qui, pour l'essentiel, reste encore à inventer, mais dont on peut esquisser les contours.

\textbf{De mémoriser à comprendre.} Si la restitution brute perd de sa valeur, la compréhension profonde des mécanismes — pourquoi une formule fonctionne, quel est le principe logique sous-jacent — reste indispensable. C'est elle qui permet de reconnaître si la réponse de l'IA est absurde ou brillante.

\textbf{De comprendre à questionner.} L'IA générative ne produit de la qualité que si on lui pose la bonne question. L'art du prompt n'est pas une technique de dactylographie ; c'est l'art de formuler une intention intellectuelle claire. Apprendre à questionner, à problématiser, à creuser une idée devient une compétence cardinale. Le bon élève n'est plus celui qui a la bonne réponse, c'est celui qui pose la question qui fait avancer.

\textbf{De questionner à vérifier.} Nous l'avons vu : l'IA hallucine, elle simplifie, elle affirme avec aplomb ce qui est faux. La compétence la plus critique du XXI\up{e} siècle pourrait bien être la \emph{vérification critique}. Savoir croiser les sources, détecter les biais, tester les limites d'une affirmation, distinguer le plausible du prouvé. C'est l'esprit scientifique appliqué au quotidien.

\textbf{De vérifier à créer.} L'IA excelle dans la recombinaison de l'existant. Elle peut imiter, synthétiser, varier. Mais la création authentique — l'idée qui émerge d'une expérience singulière, d'une intuition ou d'un point de vue inédit — reste un territoire humain. L'école doit encourager la création, non pas comme la production d'un objet parfait, mais comme l'expression d'une subjectivité en formation.

\textbf{De créer à coopérer.} Le travail de demain sera un travail hybride, où l'humain coordonne des agents IA et collabore avec d'autres humains. Apprendre à travailler en équipe, à distribuer les tâches, à intégrer des contributions hétérogènes (humaines et artificielles) est une nécessité pratique et sociale.

\textbf{De coopérer à juger.} Enfin, et peut-être surtout, l'école doit enseigner le jugement éthique et pratique. Face à un dilemme moral, face à un choix stratégique, face à une décision qui engage des vies humaines, la machine peut éclairer, mais elle ne peut pas décider à notre place, car elle n'a ni conscience, ni responsabilité. Former des citoyens capables de juger, de trancher et d'assumer, voilà l'horizon.

Ce déplacement — mémoriser → comprendre → questionner → vérifier → créer → coopérer → juger — dessine une tout autre école. Une école où l'on évalue non plus le devoir final, mais le processus : les questions posées, les vérifications effectuées, les détours intellectifs empruntés. Une école où le professeur n'est plus le distributeur du savoir, mais l'entraîneur, le mentor, le provocateur de pensée.

\textbf{Exemple concret :} plutôt que de demander à un élève de rédiger une dissertation sur la liberté, on pourrait lui demander de générer trois dissertations contradictoires avec l'IA, d'en faire la critique argumentée, d'identifier les angles morts de chacune, et de proposer une synthèse personnelle qui dépasse la machine. Le travail intellectuel se déplace de la rédaction vers l'audit et la critique. Il ne s'agit pas de naïveté. Ce déplacement suppose des enseignants formés, des classes moins chargées, et un renoncement à la commodisation de l'évaluation (les QCM, les grilles de notation standardisées). Il suppose aussi que l'on reconnaisse une vérité difficile : l'école a longtemps été un instrument de sélection basé sur la capacité à restituer un savoir. Si cette capacité perd sa valeur économique, l'école devra trouver d'autres critères de sélection, ou bien la société devra cesser d'utiliser l'école comme principal filtre social.

La question éducative est peut-être la plus urgente de toutes. Car c'est aujourd'hui que se forment les esprits qui hériteront de cette société hybride. Si nous continuons à former des humains pour être des exécutants cognitifs, nous les préparons à être remplacés. Si nous les formons pour être des juges, des créateurs et des citoyens, nous les préparons à maîtriser le nouveau pouvoir que l'IA leur confère.

Mais ce pouvoir, justement, n'est pas uniformément distribué. L'accès à l'IA, la propriété des modèles, le contrôle des infrastructures font émerger de nouvelles géométries du pouvoir. C'est ce que nous devons explorer maintenant.
